\section{Megvalósítás}
\label{ch:implementation}

\subsection{Fejlesztési környezet és eszközök}

% TODO: Használt IDE-k, eszközök, verziókezelés, csomagkezelők bemutatása

\subsection{Backend}

A backend ASP.NET Core alapokra épül, és kontrollerek, szolgáltatások, repository-k, valamint adatbázis-kezeléssel kapcsolatos komponensek közműködéséből áll. A projektben megjelenik az Entity Framework Core használata, az automatikus migrációfuttatás, a Swagger alapú API-leírás, valamint a hitelesítési és jogosultságkezelési infrastruktúra \cite{aspnetcore}.

% TODO: Részletes kifejtés:
% - a fő entitások és kapcsolatuk bemutatása
% - a repository és service rétegek indoklása
% - a jogosultságkezelés és a CRUD műveletek védelme
% - a lapozás, keresési lehetőségek és hibakezelés megoldása

\subsection{Frontend}

A frontend Vue és TypeScript alapokra épülő egyoldalas alkalmazás. A komponensalapú felépítés, a kliensoldali routing, valamint az API klienssel való integráció lehetővé teszi, hogy a felhasználói folyamatok jól elkülöníthető modulokban jelenjenek meg \cite{vue}. A felületnek különösen fontos szerepe van a receptlista, a részletes receptnézet, a keresési felület, valamint a profil- és kommentfunkciók megjelenítésében.

\subsection{API szerződés és generálás}

Az OpenAPI specifikáció külön értéket képvisel a projektben, mert egyszerre támogatja a backend dokumentálását, a kliensoldali integrációt és a következetes interfészek kialakítását. A dolgozatban itt lehet bemutatni, hogy a specifikáció hogyan kapcsolódik a generálható klienshez, és ez milyen módon csökkenti a frontend és backend közötti eltérések kockázatát \cite{openapi}.

\subsection{Fejlesztés közben felmerült problémák}

% TODO: A fejlesztés során felmerült problémák és azok megoldásai

\section{DevOps}
\label{sec:devops}

\subsection{Docker és konténeres futtatás}

A repositoryban szereplő Dockerfile-ok és Docker Compose állományok arra utalnak, hogy a projekt fejlesztéséhez és terjesztéséhez konténeres workflow is rendelkezésre áll. Ebben az alfejezetben érdemes bemutatni a fő szolgáltatások együttfutásának módját, valamint a fejlesztői és publikációs konfigurációk közötti eltéréseket \cite{docker}.

\subsection{CI/CD}

% TODO: Folyamatos integráció és szállítás bemutatása, ha van ilyen a projektben